%% toc: P2 17 17.4
\section{Proofs for Mendelson's first-order theories and the deduction theorem}
\label{sec:annex-p02-w02c-p02-017-mendelson}

Notation is that of \cref{sec:p02-w02c-p02-017-mendelson}.  A justified
deduction is written as a numbered list of steps with their reasons.  When
several justified deductions are written one after another, references in the
later ones are renumbered; a step's dependences are fixed by the steps its
justification refers to (\cref{def:p02-w02c-p02-017-dependence}), so a piece
whose references stay inside the piece and to named earlier steps has the
same dependences wherever it is placed.  When a wf of \(\mathscr L\) is read as
a formula of the corresponding signature, the variables free in it in the
sense of \cref{def:p02-w02c-p02-017-language} are exactly the members of its
\(\FV\) (\cref{sec:fol-syntax-semantics-03}): both are generated by the same
recursion, all variables of an atomic
wf, the union for \(\Rightarrow\), the same set for \(\neg\), and removal of
\(y\) for \((\forall y)\).

\begin{annexitem}{annex:p02-w02c-p02-017-identity}{thm:p02-w02c-p02-017-identity}{Identity in a first-order theory}
\begin{proof}
Let \(\mathscr B\) be a wf.  The following is a deduction from the empty set
in \(K\).
\[
\begin{array}{rll}
 1. & (\mathscr B\Rightarrow((\mathscr B\Rightarrow\mathscr B)
       \Rightarrow\mathscr B))\Rightarrow{} & \text{(A2)}\\
    & \quad((\mathscr B\Rightarrow(\mathscr B\Rightarrow\mathscr B))
       \Rightarrow(\mathscr B\Rightarrow\mathscr B)) & \\
 2. & \mathscr B\Rightarrow((\mathscr B\Rightarrow\mathscr B)
       \Rightarrow\mathscr B) & \text{(A1)}\\
 3. & (\mathscr B\Rightarrow(\mathscr B\Rightarrow\mathscr B))
       \Rightarrow(\mathscr B\Rightarrow\mathscr B) & 2,1,\ \text{MP}\\
 4. & \mathscr B\Rightarrow(\mathscr B\Rightarrow\mathscr B) & \text{(A1)}\\
 5. & \mathscr B\Rightarrow\mathscr B & 4,3,\ \text{MP}
\end{array}
\]
Step~1 is (A2) with \(\mathscr C\) taken as
\(\mathscr B\Rightarrow\mathscr B\) and \(\mathscr D\) as \(\mathscr B\).
Step~2 is (A1) with \(\mathscr C\) taken as
\(\mathscr B\Rightarrow\mathscr B\).  Step~1 is step~2 followed by
\(\Rightarrow\) and the wf of step~3, so MP gives step~3.  Step~4 is (A1) with
\(\mathscr C\) taken as \(\mathscr B\).  Step~3 is step~4 followed by
\(\Rightarrow\) and \(\mathscr B\Rightarrow\mathscr B\), so MP gives step~5.
The schemas (A1) and (A2) admit every wf of \(\mathscr L\) for their
metavariables, so the deduction exists for every \(\mathscr B\); it has no Hyp
step and no Gen step.
\end{proof}
\end{annexitem}

\begin{annexitem}{annex:p02-w02c-p02-017-free-for}{thm:p02-w02c-p02-017-free-for}{Substitution of a term free for the variable}
\begin{proof}
Induction on the number of occurrences of \(\neg\), \(\Rightarrow\) and
quantifiers in \(\mathscr B\), for all assignments at once.  Fix \(t\) and
\(x_i\), and for an assignment \(s\) write
\(v_s\defeq\sem{t}^{\mathcal M}_{s}\) and
\(s'\defeq s[x_i\mapsto v_s]\).

\emph{Atomic.}  Let \(\mathscr B\) be \(A^n_j(u_1,\ldots,u_n)\).  Every
occurrence of \(x_i\) in it is free, so \(\mathscr B(t)\) is
\(A^n_j(u_1[t/x_i],\ldots,u_n[t/x_i])\) with the term substitution of
\cref{sec:fol-syntax-semantics-04}.  By the term-substitution lemma
\cref{lem:term-substitution},
\(\sem{u_k[t/x_i]}^{\mathcal M}_{s}
=\sem{u_k}^{\mathcal M}_{s'}\) for each \(k\), and the atomic
clause of satisfaction in \cref{sec:fol-semantics-02} gives the claim.

\emph{Negation and implication.}  Let \(\mathscr B\) be \(\neg\mathscr C\) or
\(\mathscr C\Rightarrow\mathscr D\).  Every free occurrence of \(x_i\) in
\(\mathscr C\) or \(\mathscr D\) is a free occurrence in \(\mathscr B\), and
every quantifier of \(\mathscr C\) or \(\mathscr D\) is a quantifier of
\(\mathscr B\) with the same scope.  So \(t\) is free for \(x_i\) in
\(\mathscr C\) and in \(\mathscr D\).  Since \(\mathscr B(t)\) is
\(\neg\mathscr C(t)\) or \(\mathscr C(t)\Rightarrow\mathscr D(t)\), the
induction hypothesis and the clauses for \(\neg\) and \(\to\) give the claim.

\emph{Quantifier.}  Let \(\mathscr B\) be \((\forall y)\mathscr C\).

Suppose first that \(x_i\) is not free in \(\mathscr B\); this includes the
case \(y=x_i\).  Then \(\mathscr B(t)\) is \(\mathscr B\), and \(s\) and
\(s'\) differ at most at \(x_i\notin\FV(\mathscr B)\), so formula coincidence
\cref{lem:fol-coincidence} gives
\(\mathcal M,s\entails\mathscr B\) iff \(\mathcal M,s'\entails\mathscr B\).

Suppose now that \(x_i\) is free in \(\mathscr B\).  Then \(y\neq x_i\) and
\(x_i\) has a free occurrence in \(\mathscr C\), which is a free occurrence in
\(\mathscr B\) within the scope of \((\forall y)\).  As \(t\) is free for
\(x_i\) in \(\mathscr B\), \(y\) is not a variable of \(t\).  Every free
occurrence of \(x_i\) in \(\mathscr C\) is free in \(\mathscr B\) because
\(y\neq x_i\), and the quantifiers of \(\mathscr C\) are quantifiers of
\(\mathscr B\) with the same scopes, so \(t\) is free for \(x_i\) in
\(\mathscr C\).  Also \(\mathscr B(t)\) is \((\forall y)\mathscr C(t)\).  For
\(a\in\dom\mathcal M\) put \(s_a\defeq s[y\mapsto a]\).  By the term-coincidence lemma of
\cref{sec:fol-semantics-01}, \(v_{s_a}=v_s\), since \(s_a\) and \(s\)
agree on every variable of \(t\).  The induction hypothesis for \(\mathscr C\)
at \(s_a\) gives
\(\mathcal M,s_a\entails\mathscr C(t)\) iff
\(\mathcal M,s_a[x_i\mapsto v_s]\entails\mathscr C\), and
\(s_a[x_i\mapsto v_s]=s'[y\mapsto a]\) because \(y\neq x_i\).  Hence
\[
 \mathcal M,s\entails(\forall y)\mathscr C(t)
 \iff\text{for all }a,\ \mathcal M,s'[y\mapsto a]\entails\mathscr C
 \iff\mathcal M,s'\entails(\forall y)\mathscr C.
\]
This completes the induction.
\end{proof}
\end{annexitem}

\begin{annexitem}{annex:p02-w02c-p02-017-soundness}{thm:p02-w02c-p02-017-soundness}{Theorems of a first-order theory hold in its models}
\begin{proof}
Let \(\mathscr D_1,\ldots,\mathscr D_n\) be a deduction from the empty set in
\(K\).  We show by strong induction on \(i\) that \(\mathcal M,s\entails
\mathscr D_i\) for every assignment \(s\).  With no hypotheses, each step is
an axiom or follows by MP or Gen.  By the clauses for \(\neg\) and \(\to\) in
\cref{sec:fol-semantics-02}, whether \(s\) satisfies a wf built by
\(\neg\) and \(\Rightarrow\) from given wfs depends only on whether \(s\)
satisfies those wfs, by the truth tables of negation and material
implication.

\emph{(A1).}  If \(s\) satisfies \(\mathscr B\), it satisfies
\(\mathscr C\Rightarrow\mathscr B\); so it satisfies
\(\mathscr B\Rightarrow(\mathscr C\Rightarrow\mathscr B)\).

\emph{(A2).}  Let \(s\) satisfy
\(\mathscr B\Rightarrow(\mathscr C\Rightarrow\mathscr D)\) and
\(\mathscr B\Rightarrow\mathscr C\).  If \(s\) satisfies \(\mathscr B\), it
satisfies \(\mathscr C\Rightarrow\mathscr D\) and \(\mathscr C\), hence
\(\mathscr D\).  So \(s\) satisfies \(\mathscr B\Rightarrow\mathscr D\), and
therefore every instance of (A2).

\emph{(A3).}  Let \(s\) satisfy \(\neg\mathscr C\Rightarrow\neg\mathscr B\)
and \(\neg\mathscr C\Rightarrow\mathscr B\).  If \(s\) did not satisfy
\(\mathscr C\), it would satisfy \(\neg\mathscr C\), hence both
\(\neg\mathscr B\) and \(\mathscr B\), which is impossible.  So \(s\)
satisfies \(\mathscr C\), and therefore every instance of (A3).

\emph{(A4).}  Let \(t\) be free for \(x_i\) in \(\mathscr B(x_i)\) and let
\(s\) satisfy \((\forall x_i)\mathscr B(x_i)\).  By the quantifier clause,
\(s[x_i\mapsto\sem{t}^{\mathcal M}_{s}]\) satisfies
\(\mathscr B(x_i)\), so \(s\) satisfies \(\mathscr B(t)\) by
\cref{thm:p02-w02c-p02-017-free-for}.

\emph{(A5).}  Let \(x_i\) not be free in \(\mathscr B\), and let \(s\) satisfy
\((\forall x_i)(\mathscr B\Rightarrow\mathscr C)\) and \(\mathscr B\).  For
\(a\in\dom\mathcal M\), \(s[x_i\mapsto a]\) agrees with \(s\) on
\(\FV(\mathscr B)\), so it satisfies \(\mathscr B\) by
\cref{lem:fol-coincidence}; it also satisfies
\(\mathscr B\Rightarrow\mathscr C\), hence \(\mathscr C\).  So \(s\)
satisfies \((\forall x_i)\mathscr C\), and therefore every instance of (A5).

\emph{Proper axioms} are satisfied by every assignment by assumption.

\emph{MP.}  If every assignment satisfies \(\mathscr D_j\) and
\(\mathscr D_j\Rightarrow\mathscr D_i\), every assignment satisfies
\(\mathscr D_i\).

\emph{Gen.}  Let \(\mathscr D_i\) be \((\forall x_k)\mathscr D_j\) with
\(j<i\), and let every assignment satisfy \(\mathscr D_j\).  For an
assignment \(s\) and \(a\in\dom\mathcal M\), \(s[x_k\mapsto a]\) is an
assignment and satisfies \(\mathscr D_j\); so \(s\) satisfies
\(\mathscr D_i\).

The last step of a proof of a theorem is therefore satisfied by every
assignment.  A predicate calculus has no proper axioms, so the assumption on
\(\mathcal M\) is empty and the conclusion holds in every structure.
\end{proof}
\end{annexitem}

\begin{annexitem}{annex:p02-w02c-p02-017-prop24}{thm:p02-w02c-p02-017-prop24}{Mendelson's Proposition 2.4}
\begin{proof}
Let \(\mathscr D_1,\ldots,\mathscr D_n\) be a justified deduction from
\(\Gamma\cup\{\mathscr B\}\) with \(\mathscr D_n=\mathscr C\), and let
\(\mathscr C\) not depend upon \(\mathscr B\) in it.  Let \(I\) be the set of
indices \(i\) such that \(\mathscr D_i\) does not depend upon \(\mathscr B\);
\(n\in I\).  Let \(S\) be the subsequence of the steps with index in \(I\), in
their order, each keeping its justification.  We check that \(S\) is a
justified deduction from \(\Gamma\).  Let \(i\in I\).
\begin{itemize}
  \item If \(\mathscr D_i\) is justified by Hyp, then \(\mathscr D_i\) is not
        \(\mathscr B\), for otherwise it would depend upon \(\mathscr B\) by
        clause~1 of \cref{def:p02-w02c-p02-017-dependence}.  So
        \(\mathscr D_i\in\Gamma\).
  \item If \(\mathscr D_i\) is an axiom, it remains one.
  \item If \(\mathscr D_i\) follows by MP from \(\mathscr D_j\) and
        \(\mathscr D_k\), or by Gen from \(\mathscr D_j\), then none of these
        premisses depends upon \(\mathscr B\), for otherwise \(\mathscr D_i\)
        would by clause~2.  So their indices lie in \(I\), they precede
        \(\mathscr D_i\) in \(S\), and the same rule applies.
\end{itemize}
\(S\) ends with \(\mathscr D_n=\mathscr C\), so \(\Gamma\vdash_K\mathscr C\).

For later use we note two properties of \(S\).  (P1) Every Gen step of
\(S\) is a Gen step of the given deduction, with the same premiss and the same
quantified variable.  (P2) For every \(\mathscr E\in\Gamma\) and
\(i\in I\), \(\mathscr D_i\) depends upon \(\mathscr E\) in \(S\) if and only
if it depends upon \(\mathscr E\) in the given deduction.  (P1) holds because
\(S\) keeps the justifications.  (P2) follows by induction along \(I\): a Hyp
step depends upon \(\mathscr E\) in either deduction exactly when it is
\(\mathscr E\); an axiom step depends upon nothing in either; an MP or Gen
step refers to the same premisses in both, which by the induction hypothesis
have the same dependences.

Mendelson argues by induction on the length of the deduction
\cite[\S 2.4, p.~74]{Mendelson1997}; the subsequence \(S\) is the deduction
that induction produces.
\end{proof}
\end{annexitem}

\begin{annexitem}{annex:p02-w02c-p02-017-prop25}{thm:p02-w02c-p02-017-prop25}{Mendelson's Proposition 2.5, the deduction theorem}
\begin{proof}
Let \(\mathscr D_1,\ldots,\mathscr D_n\) be a justified deduction from
\(\Gamma\cup\{\mathscr B\}\) with \(\mathscr D_n=\mathscr C\), in which no
application of Gen to a wf depending upon \(\mathscr B\) has as quantified
variable a free variable of \(\mathscr B\).  We build, for
\(i=1,\ldots,n\) in turn, a piece \(\mathcal E_i\) of steps so that the
concatenation \(\mathcal E_1\cdots\mathcal E_i\) is a justified deduction from
\(\Gamma\) whose last step is \(\mathscr B\Rightarrow\mathscr D_i\), and so
that
\begin{itemize}
  \item[(Q)] for every \(\mathscr E\in\Gamma\), if the last step of
        \(\mathcal E_i\) depends upon \(\mathscr E\), then \(\mathscr D_i\)
        depends upon \(\mathscr E\) in the given deduction.
\end{itemize}
Each piece refers only to its own steps and to last steps of earlier pieces.
The cases follow the justification of \(\mathscr D_i\); \(j,k<i\).

\emph{Case 1: \(\mathscr D_i\) is justified by Hyp and is \(\mathscr B\).}
\(\mathcal E_i\) is the five-step deduction of
\(\mathscr B\Rightarrow\mathscr B\) of
\cref{thm:p02-w02c-p02-017-identity}.  It has no Hyp step, so its last step
depends upon nothing and (Q) holds.

\emph{Case 2: \(\mathscr D_i\) is justified by Hyp and is not
\(\mathscr B\).}  Then \(\mathscr D_i\in\Gamma\), and \(\mathcal E_i\) is
\[
 \mathscr D_i\ \text{(Hyp)},\qquad
 \mathscr D_i\Rightarrow(\mathscr B\Rightarrow\mathscr D_i)\ \text{(A1)},\qquad
 \mathscr B\Rightarrow\mathscr D_i\ \text{(MP)}.
\]
Its last step depends upon \(\mathscr D_i\) and upon no other member of
\(\Gamma\), and \(\mathscr D_i\) depends upon itself in the given deduction
by clause~1.

\emph{Case 3: \(\mathscr D_i\) is an axiom.}  \(\mathcal E_i\) is
\(\mathscr D_i\) (axiom),
\(\mathscr D_i\Rightarrow(\mathscr B\Rightarrow\mathscr D_i)\) (A1),
\(\mathscr B\Rightarrow\mathscr D_i\) (MP).  Its last step depends upon
nothing.

\emph{Case 4: \(\mathscr D_i\) follows by MP from \(\mathscr D_j\) and
\(\mathscr D_k=(\mathscr D_j\Rightarrow\mathscr D_i)\).}  The last steps of
\(\mathcal E_j\) and \(\mathcal E_k\) are \(\mathscr B\Rightarrow\mathscr D_j\)
and \(\mathscr B\Rightarrow(\mathscr D_j\Rightarrow\mathscr D_i)\).
\(\mathcal E_i\) is
\[
\begin{array}{ll}
 (\mathscr B\Rightarrow(\mathscr D_j\Rightarrow\mathscr D_i))\Rightarrow
 ((\mathscr B\Rightarrow\mathscr D_j)\Rightarrow
  (\mathscr B\Rightarrow\mathscr D_i)) & \text{(A2)}\\
 (\mathscr B\Rightarrow\mathscr D_j)\Rightarrow
  (\mathscr B\Rightarrow\mathscr D_i)
   & \text{MP, last step of }\mathcal E_k\\
 \mathscr B\Rightarrow\mathscr D_i
   & \text{MP, last step of }\mathcal E_j.
\end{array}
\]
If its last step depends upon \(\mathscr E\), so does the last step of
\(\mathcal E_j\) or of \(\mathcal E_k\); by (Q) for \(j\) and \(k\),
\(\mathscr D_j\) or \(\mathscr D_k\) depends upon \(\mathscr E\), and then
\(\mathscr D_i\) does by clause~2.

\emph{Case 5: \(\mathscr D_i\) is \((\forall x_m)\mathscr D_j\), by Gen with
quantified variable \(x_m\).}

\emph{Case 5a: \(\mathscr D_j\) does not depend upon \(\mathscr B\).}  The
initial segment \(\mathscr D_1,\ldots,\mathscr D_j\) is a justified deduction
from \(\Gamma\cup\{\mathscr B\}\) in which \(\mathscr D_j\) does not depend
upon \(\mathscr B\) (\cref{def:p02-w02c-p02-017-dependence}).  Let \(S_j\) be
the deduction of \(\mathscr D_j\) from \(\Gamma\) built from that segment in
the proof of \cref{thm:p02-w02c-p02-017-prop24}.  \(\mathcal E_i\) is
the steps of \(S_j\) followed by
\[
\begin{array}{ll}
 (\forall x_m)\mathscr D_j & \text{Gen from the last step of }S_j\\
 \mathscr D_i\Rightarrow(\mathscr B\Rightarrow\mathscr D_i) & \text{(A1)}\\
 \mathscr B\Rightarrow\mathscr D_i & \text{MP}.
\end{array}
\]
The Hyp steps of \(S_j\) are members of \(\Gamma\).  If the last step of
\(\mathcal E_i\) depends upon \(\mathscr E\), so does the Gen step, hence the
last step of \(S_j\), hence, by (P2) of that proof, \(\mathscr D_j\) in the
given deduction, and then \(\mathscr D_i\) by clause~2.

\emph{Case 5b: \(\mathscr D_j\) depends upon \(\mathscr B\).}  Step~\(i\)
applies Gen to a wf depending upon \(\mathscr B\), so by hypothesis \(x_m\)
is not free in \(\mathscr B\), and
\((\forall x_m)(\mathscr B\Rightarrow\mathscr D_j)\Rightarrow
(\mathscr B\Rightarrow(\forall x_m)\mathscr D_j)\) is an instance of (A5).
\(\mathcal E_i\) is
\[
\begin{array}{ll}
 (\forall x_m)(\mathscr B\Rightarrow\mathscr D_j)
   & \text{Gen from the last step of }\mathcal E_j\\
 (\forall x_m)(\mathscr B\Rightarrow\mathscr D_j)\Rightarrow
 (\mathscr B\Rightarrow(\forall x_m)\mathscr D_j) & \text{(A5)}\\
 \mathscr B\Rightarrow\mathscr D_i & \text{MP}.
\end{array}
\]
If its last step depends upon \(\mathscr E\), so does the Gen step, hence the
last step of \(\mathcal E_j\); by (Q) for \(j\), \(\mathscr D_j\) depends upon
\(\mathscr E\), and so does \(\mathscr D_i\).

In every case the Hyp steps of \(\mathcal E_i\) are members of \(\Gamma\),
and each step is an axiom or follows by MP or Gen from earlier steps, so
\(\mathcal E_1\cdots\mathcal E_i\) is a justified deduction from \(\Gamma\).
For \(i=n\) its last step is \(\mathscr B\Rightarrow\mathscr C\), so
\(\Gamma\vdash_K\mathscr B\Rightarrow\mathscr C\).

Mendelson's proof has the same case division; in Case~1 he cites his
Proposition~2.1 for \(\vdash\mathscr B\Rightarrow\mathscr B\)
\cite[\S 2.4, p.~74]{Mendelson1997}, which is replaced here by
\cref{thm:p02-w02c-p02-017-identity}.
\end{proof}
\end{annexitem}

\begin{annexitem}{annex:p02-w02c-p02-017-cor26}{thm:p02-w02c-p02-017-cor26}{Mendelson's Corollary 2.6}
\begin{proof}
Let the given deduction of \(\mathscr C\) from \(\Gamma\cup\{\mathscr B\}\),
with any justification, contain no Gen step whose quantified variable is free
in \(\mathscr B\).  Then in particular no Gen step applied to a wf depending
upon \(\mathscr B\) has such a variable, which is the hypothesis of
\cref{thm:p02-w02c-p02-017-prop25} for this deduction.  That proposition gives
\(\Gamma\vdash_K\mathscr B\Rightarrow\mathscr C\).
\end{proof}
\end{annexitem}

\begin{annexitem}{annex:p02-w02c-p02-017-cor27}{thm:p02-w02c-p02-017-cor27}{Mendelson's Corollary 2.7}
\begin{proof}
Since \(\Gamma,\mathscr B\vdash_K\mathscr C\), there is a deduction of
\(\mathscr C\) from \(\Gamma\cup\{\mathscr B\}\); give each step one of the
reasons that make it a deduction step.  As \(\mathscr B\) is closed, no
variable is free in \(\mathscr B\), so no Gen step of this deduction has a
quantified variable free in \(\mathscr B\).
\Cref{thm:p02-w02c-p02-017-cor26} gives
\(\Gamma\vdash_K\mathscr B\Rightarrow\mathscr C\).
\end{proof}
\end{annexitem}

\begin{annexitem}{annex:p02-w02c-p02-017-extension}{thm:p02-w02c-p02-017-extension}{Extension of Propositions 2.4--2.7}
\begin{proof}
\emph{Proposition 2.4.}  Take the deduction \(S\) of
Annex~\ref{annex:p02-w02c-p02-017-prop24}.  By (P1) a Gen step of \(S\) is a Gen
step of the given deduction with the same quantified variable and the same
premiss step; by (P2) that premiss depends upon \(\mathscr E\in\Gamma\) in
\(S\) only if it does in the given deduction.

\emph{Proposition 2.5.}  Take the deduction \(\mathcal E_1\cdots\mathcal E_n\)
of Annex~\ref{annex:p02-w02c-p02-017-prop25}.  Gen steps occur only in pieces of
Cases~5a and~5b, where \(\mathscr D_i\) is \((\forall x_m)\mathscr D_j\) by
Gen in the given deduction.
\begin{itemize}
  \item In Case~5a, the Gen steps inside \(S_j\) are, by (P1) and (P2) applied
        to the segment \(\mathscr D_1,\ldots,\mathscr D_j\), Gen steps of that
        segment with the same variables and premisses whose dependences on
        members of \(\Gamma\) are those of the segment, which are those of the
        whole deduction.  The further Gen step, with variable \(x_m\), is
        applied to the last step of \(S_j\), which depends upon
        \(\mathscr E\) only if \(\mathscr D_j\) does, by (P2); and step~\(i\)
        of the given deduction applies Gen with variable \(x_m\) to
        \(\mathscr D_j\).
  \item In Case~5b, the Gen step has variable \(x_m\) and is applied to the
        last step of \(\mathcal E_j\), which depends upon \(\mathscr E\) only
        if \(\mathscr D_j\) does, by (Q); and step~\(i\) of the given
        deduction applies Gen with variable \(x_m\) to \(\mathscr D_j\).
\end{itemize}
In each case a Gen step of the new deduction applied to a wf depending upon
\(\mathscr E\) is matched by a Gen step of the given deduction with the same
variable applied to a wf depending upon \(\mathscr E\).

\emph{Corollaries 2.6 and 2.7.}  Their proofs apply
\cref{thm:p02-w02c-p02-017-prop25} to the given deduction itself, so the
deduction \(\mathcal E_1\cdots\mathcal E_n\) built from it has the property
just shown.
\end{proof}
\end{annexitem}

\begin{annexitem}{annex:p02-w02c-p02-017-failure}{thm:p02-w02c-p02-017-failure}{The unmodified deduction theorem fails}
\begin{proof}
\emph{Derivability.}  The sequence \(A^1_1(x_1)\) (Hyp),
\((\forall x_1)A^1_1(x_1)\) (Gen from step~1, quantified variable \(x_1\)) is
a deduction from \(\{A^1_1(x_1)\}\), since Gen has no side condition
(\cref{def:p02-w02c-p02-017-theory}).

\emph{Non-theoremhood.}  Let \(\mathcal M\) have domain \(\{0,1\}\),
interpret \(A^1_1\) as \(\{0\}\), and interpret the remaining symbols of the
language arbitrarily (every function and constant symbol by the value \(0\),
every other relation symbol by the empty relation).  Let \(s\) be an
assignment with \(s(x_1)=0\).  Then \(\mathcal M,s\entails A^1_1(x_1)\), and
\(s[x_1\mapsto1]\) does not satisfy \(A^1_1(x_1)\), so
\(\mathcal M,s\not\entails(\forall x_1)A^1_1(x_1)\).  By the clause for
\(\to\), \(s\) does not satisfy
\(A^1_1(x_1)\Rightarrow(\forall x_1)A^1_1(x_1)\).  \(K\) is a predicate
calculus, so by \cref{thm:p02-w02c-p02-017-soundness} every theorem of \(K\)
is satisfied by \(s\) in \(\mathcal M\); hence this wf is not a theorem of
\(K\).

\emph{The hypothesis is violated.}  If some justified deduction of
\((\forall x_1)A^1_1(x_1)\) from \(\{A^1_1(x_1)\}\) satisfied the hypothesis of
\cref{thm:p02-w02c-p02-017-prop25} with \(\Gamma=\varnothing\) and
\(\mathscr B=A^1_1(x_1)\), that proposition would give
\(\vdash_K A^1_1(x_1)\Rightarrow(\forall x_1)A^1_1(x_1)\), contrary to what
was just shown.  In the two-step deduction above the violation is visible:
step~2 applies Gen with variable \(x_1\), free in \(A^1_1(x_1)\), to step~1,
which depends upon \(A^1_1(x_1)\).
\end{proof}
\end{annexitem}
